% SPDX-License-Identifier: CC-BY-4.0
% Copyright (c) 2026 Martin Schröder <info@swedishembedded.com>

\section{What the loop needs in order to work}\label{sec:requirements}

A turn of the loop adds a capability only if its stages cannot be satisfied by an artefact of
the setup. This section states the requirements we found necessary, the mechanism that
implements each, and which are not yet implemented. Each names its commit or file in the
ledger, which checks that the commit is in the history of the pinned revision of its
repository (\cref{sec:repro}).

\paragraph{Held-out units must be independent.} Records carry the group of the source text they
are grounded in, and a held-out measurement needs whole groups. A dataset that is one source
text cannot be divided, so the split refuses it with a typed error naming the cause instead of
dividing the group and failing the leakage check afterwards; a free record follows the fixed
side of its group, monitoring with no whole group to set aside is skipped with the reason
recorded, and a source that is one independent text is refused at intake before any rollout is
paid for.\cl{FX01,FX05} A source of independent units is therefore one file per unit.

\paragraph{Admission must see through markup.} The subject of a task is compared with the
instruction, the source and the variants through one function that ignores code quotes,
quotation marks, emphasis and a call's empty parentheses, so that ``the \code{q} key'' contains
the subject ``q key''.\cl{FX02} A number written with a thousands separator is the same number
in the stated-reference check.\cl{FX06}

\paragraph{The serve check must start the real server.} The gate's serve check runs
\code{brain serve --adapter}. A candidate's card id must follow Brain's adapter grammar
(base, owner, name, tag), and the check must pass the bank option that covers every linear
Splinter's LoRA corrects; the test double of the server refuses an adapter offered without it,
because a double that accepts any flags hides both.\cl{FX03,FX04}

\paragraph{The anchor bound must count net lost answers on a large enough suite.} A rate bound of
0.02 on a 40-question anchor suite allows no lost answer, since one net loss is 0.025. The gate
can decide the anchor on answers lost net of answers gained (an option; the default rate bound
is unchanged), and a 200-question anchor suite built by code and a reviewed list
exists.\cl{FX07}

\paragraph{Retention must test acquired capabilities.} The retention suite of each earlier
release is built from that release's held-out split, and the exams of workspace families are
added to it.\cl{OP01} For a domain of facts the held-out split holds facts the release never
taught, so a passed retention check is not evidence that taught facts were kept. A retention
exam drawn from what each release demonstrably acquired is planned and not implemented; the
registered coding-task retention check (\cref{sec:protocol}) uses the exam families of every
earlier round for this reason.

\paragraph{The verifier must prove that the tests ran.} The workspace verifier passes a patch
only if every fail-to-pass and pass-to-pass command exits with status 0.\cl{OP02} That does not
prove that the intended tests ran, so a suite that exits zero without running them would pass.
A check that the named tests were executed and passed is planned; the coding-task experiments
register conformance measurements for the verifier before they count (\cref{sec:protocol}).

\paragraph{The gate must carry canaries.} A gate with an empty canary set skips that check. The
release gate called from \code{practice} is built without canary tasks,\cl{OP03} so a release
from \code{practice} does not yet carry that evidence; adding canaries is planned.

\subsection{Trajectory and training requirements of repository tasks}\label{sec:audit}

Training an agent on its own repository trajectories adds requirements to the stack. Each is
located in source; the ledger names the file, commit or test.

\begin{enumerate}
\item \textbf{Loss masking of tool conversations} (Brain; implemented). Under the Qwen3 template
an empty reasoning block is added to the last assistant turn only, so non-final assistant
turns of a multi-step conversation failed the trainer's boundary check. Tool-calling
conversations now train per assistant turn in both thinking modes.\cl{AU01}
\item \textbf{Tool definitions lost between recording and training} (implemented). The agent
records the tool definitions of the run, but the trajectory projections dropped them; an
ATIF trajectory's tool definitions now become the chat sample's tools, and Splinter's
dataset check covers tools.\cl{AU02}
\item \textbf{Trajectory fidelity in the agent SDK} (Sven; implemented). The exported trajectory
marked no context compaction and kept the uncompacted history, listed every registered tool
rather than those offered, and turned unparseable tool-call arguments into an empty object.
Compactions are now marked and the compacted history kept, the offered tools are exported,
and the raw text of non-JSON arguments is kept.\cl{AU03}
\item \textbf{Release probes were closed-book} (implemented). A coding agent is now gated on the
exam its run reserved, with report types for workspace suites.\cl{AU04}
\item \textbf{Contamination paths} (implemented). The semantic-memory tool is withheld from the
worker, and a context tool that could open any readable path is jailed.\cl{AU05}
\item \textbf{Silent waits} (Sven; implemented). A model server that goes quiet is bounded and
the wait is reported, with a stalling mock provider to test it.\cl{AU06}
\item \textbf{Tool-definition mismatches} (Sven; implemented). Definitions disagreed with the
implementation and were tightened.\cl{AU07}
\item \textbf{Serving several adapters} (Brain; implemented in the qwen3 server). An adapter bank
lets one batch serve the base and many LoRA adapters.\cl{AU08}
\item \textbf{Containment} (planned). A coding agent's shell is not a sandbox. Where no
container is configured every environment record says \code{isolation: none}, and the
exposure audit detects exposure but cannot prevent it.\cl{AU09}
\end{enumerate}
